\section*{अध्याय \ref{ch:b1:lipids} --- लिपिड}

\begin{solution}{exo:b1:lipids:1}
कोई \omterm{def:b1:lipids:fattyacid}{लिपिड} ऐसा जैविक अणु है जो जल में अविलेय और अध्रुवी विलायकों में विलेय
है। \omterm{def:b1:lipids:triglyceride}{वसाएँ}, \omterm{def:b1:lipids:amphiphilic}{फ़ॉस्फ़ोलिपिड}, \omterm{def:b1:lipids:amphiphilic}{स्टेरॉल} और कैरोटिनॉइड कोई साझा
कंकाल नहीं बाँटते, केवल जल के प्रति यही बरताव, और यही उन्हें उनकी साझी
भूमिकाएँ देता है (भंडार, फ़िल्म, झिल्लियाँ)।
\end{solution}

\begin{solution}{exo:b1:lipids:2}
C18:2 $\Delta^{9,12}$। आखिरी द्विबंध कार्बॉक्सिल से गिनते हुए कार्बन 12
पर शुरू होता है, यानी मेथिल सिरे से कार्बन 6 पर: ओमेगा-6.
\end{solution}

\begin{solution}{exo:b1:lipids:3}
किसी \omterm{def:b1:lipids:triglyceride}{ट्राइग्लिसराइड} का कोई ध्रुवी शीर्ष नहीं होता: पूरी तरह जलविरोधी
होने से वह जल से एक थोक प्रावस्था के रूप में, यानी एक बूँद के रूप में बाहर
कर दिया जाता है। कोई \omterm{def:b1:lipids:amphiphilic}{फ़ॉस्फ़ोलिपिड} \omterm{def:b1:lipids:amphiphilic}{उभयरागी} है: उसके ध्रुवी शीर्ष को जल में
रहना ही है जबकि उसकी पूँछों को उसे छोड़ना ही है, और यह केवल दो अणु मोटी
कोई चादर ही पूरा कर सकती है।
\end{solution}

\begin{solution}{exo:b1:lipids:4}
लगभग \qty{41}{\celsius}, \qty{-5}{\celsius} और \qty{20}{\celsius}
(\omterm{def:b1:lipids:amphiphilic}{कोलेस्टेरॉल} के वक्र में कोई तीखा संक्रमण नहीं है; उसकी आधी चढ़ाई
\qty{20}{\celsius} के पास है)। \qty{10}{\celsius} पर किसी बैक्टीरियम को
\omterm{def:b1:lipids:fattyacid}{असंतृप्त} संघटन चाहिए, जो उसके वृद्धि-ताप से बहुत नीचे तक तरल रहे।
\end{solution}

\begin{solution}{exo:b1:lipids:5}
$15\,000\times 38 = \qty{570}{MJ}$; $570/8 = 71$ दिन। ग्लाइकोजन: शुष्क
$570\,000/17 = \qty{33.5}{kg}$, जलयोजित \qty{134}{kg}।
\end{solution}

\begin{solution}{exo:b1:lipids:6}
C22:0 $>$ C18:0 $>$ C14:0 $>$ C18:1 $>$ C18:3. संतृप्त अम्लों में
लंबी शृंखलाएँ अधिक \omterm{def:b1:water-small-molecules:bonds}{वान डर वाल्स} संपर्क बनाती हैं; एक \emph{सिस} द्विबंध
जितना जमाव हटाता है उतना चार अतिरिक्त कार्बन नहीं जोड़ते (C18:1 का गलनांक
C14:0 से नीचे है); और हर आगे का बंध उसे फिर नीचे ले जाता है।
\end{solution}

\begin{solution}{exo:b1:lipids:7}
$2\times\qty{140}{\micro m^2}/\qty{0.6}{nm^2} = 2\times 1.4\times
10^{-10}/6\times 10^{-19} = \num{4.7e8}$ अणु; द्रव्यमान
$\num{4.7e8}\times 750/\num{6e23} = \qty{5.8e-13}{g} = \qty{0.58}{pg}$।
\end{solution}

\begin{solution}{exo:b1:lipids:8}
आयतन \qty{1.11}{mL} $= \qty{1.11e-6}{m^3}$; गोलों का पृष्ठ
$= 6V/d = 6\times 1.11\times 10^{-6}/10^{-6} = \qty{6.7}{m^2}$। उतने ही
आयतन की एक बूँद: $d = \qty{1.28}{cm}$, पृष्ठ \qty{5.2}{cm^2}: यानी
पायस में तेरह हज़ार गुना अधिक अंतरापृष्ठ है। लाइपेस जल में विलेय है
और केवल अंतरापृष्ठ पर काम करती है; दर उसी के अनुपात में है।
\end{solution}

\begin{solution}{exo:b1:lipids:9}
$T_m$ से ऊपर उसके कठोर वलय पड़ोसी शृंखलाओं की गति में बाधा डालते हैं और
\omterm{def:b1:lipids:fluidity}{तरलता} घटाते हैं; $T_m$ से नीचे उसका भारीपन शृंखलाओं को क्रमबद्ध जेल में
जमने नहीं देता और जमने से रोकता है। कोई सहकारी जमाव संभव न होने से ऐसा
कोई तापमान नहीं रहता जिस पर समूचा द्विस्तर एक साथ अवस्था बदले: संक्रमण
फैल जाता है।
\end{solution}

\begin{solution}{exo:b1:lipids:10}
सब संतृप्त शृंखलाएँ: \qty{25}{\celsius} पर झिल्लियाँ पहले ही
अपने संक्रमण के पास और अकड़ी हैं; \qty{5}{\celsius} पर वे जेल बन जाती हैं।
जेल बनी झिल्लियाँ प्रोटीनों को काम करना बंद करा देती हैं और, फिर गरम होने
पर या यांत्रिक प्रतिबल पर, चटककर रिसती हैं: \omterm{def:b1:cell-unit-of-life:cell}{कोशिकाएँ} अपनी सामग्री खो देती
हैं और \omterm{def:b1:body-plans-tissues:tissue}{ऊतक} मर जाता है (शीत-क्षति)। वन्य प्ररूप शरद में अपने
\omterm{def:b1:lipids:fattyacid}{लिपिड} \omterm{def:b1:lipids:fattyacid}{असंतृप्त} कर लेता है और \qty{5}{\celsius} पर तरल बना रहता है।
\end{solution}

\begin{solution}{exo:b1:lipids:11}
प्रति किलोग्राम \omterm{def:b1:lipids:triglyceride}{वसा} \qty{1.07}{kg} जल। ऑक्सीजन: \qty{2000}{L};
\qty{21}{\%} ऑक्सीजन के \qty{5}{\%} निष्कर्षण पर, चाहिए हवा
$2000/(0.21\times 0.05) = \qty{190}{m^3}$; और \qty{30}{g/m^3} पर खोया
जल: \qty{5.7}{kg}। ऊँट साँस लेने में \omterm{def:b1:lipids:triglyceride}{वसा} से मिलने वाले जल से पाँच गुना
अधिक जल खोता है: कूबड़ ऊर्जा का भंडार है, जल का
भंडार नहीं।
\end{solution}

\begin{solution}{exo:b1:lipids:12}
कोई द्विस्तर दे दिया जाए तो जलविरोधी प्रभाव उसे फिर से बंद करा देता है,
पुटिकाओं में बँधवा देता है और बिना किसी ऊर्जा-निवेश के नए \omterm{def:b1:lipids:fattyacid}{लिपिड} ग्रहण
करवा देता है: यानी जुड़ना स्वतःस्फूर्त है। पर \omterm{def:b1:cell-unit-of-life:cell}{कोशिका} अपने \omterm{def:b1:lipids:fattyacid}{लिपिड} ऐसे
एंजाइमों से बनाती है जो किसी पहले से बनी झिल्ली (ER) में बैठे हैं, और वे
उन्हें वहीं धँसा देते हैं; कोई नई झिल्ली किसी पुरानी के फैलाव से बढ़ती है
और विभाजन पर बाँट दी जाती है, ताकि हर \omterm{def:b1:cell-unit-of-life:cell}{कोशिका} की हर झिल्ली पिछली \omterm{def:b1:cell-unit-of-life:cell}{कोशिका} की
झिल्लियों से उतरे। स्वतः-संयोजन चादर की भौतिकी समझाता है, \omterm{def:b1:cell-unit-of-life:cell}{कोशिका} में उसका
उद्गम नहीं।
\end{solution}

\begin{solution}{pb:b1:lipids:1}
\textbf{1.} C16:0 (0), C18:1 $\Delta^9$ (1), C22:6 $\Delta^{4,7,10,13,16,19}$
(6)।
\textbf{2.} गरम: $0.40\times 0 + 0.45\times 1 + 0.15\times 6 = 1.35$;
ठंडा: $0 + 0.40 + 2.10 = 2.50$ द्विबंध प्रति शृंखला।
\textbf{3.} हर \emph{सिस} बंध शृंखला में मोड़ डालता है और पड़ोसियों के
साथ कसकर जमने नहीं देता; कम \omterm{def:b1:water-small-molecules:bonds}{वान डर वाल्स} संपर्क, शृंखलाओं को अव्यवस्थित
करने के लिए कम ऊष्मा, और नीचा $T_m$।
\textbf{4.} प्रति शृंखला $1.15$ अधिक बंध: लगभग \qty{45}{\celsius}
नीचे।
\textbf{5.} किसी जेल में \omterm{def:b1:lipids:fattyacid}{लिपिड} चल ही नहीं सकते: वाहक अपनी संरूपण नहीं
बदल सकते, पंप रुक जाते हैं, \omterm{def:b1:membranes-transport:transporters}{चैनल} अटक जाते हैं; आयन-प्रवणताएँ क्षीण हो
जाती हैं, तंत्रिका-चालन विफल हो जाता है, और मछली उस ठंड से लकवाग्रस्त
होकर मर जाती है जिसे कोई ठंड-अभ्यस्त मछली सह लेती है।
\textbf{6.} असंतृप्तक, जो वसा-ऐसिल शृंखलाओं में द्विबंध डालते हैं; वे
\omterm{def:b1:eukaryotic-cell:endomembrane}{अंतर्द्रव्यी जालिका} के समाकल प्रोटीन हैं, जहाँ
\omterm{def:b1:lipids:fattyacid}{लिपिड} बनते हैं।
\textbf{7.} दोनों में वह \omterm{def:b1:lipids:fluidity}{तरलता} को सँभालता है: बहुत \omterm{def:b1:lipids:fattyacid}{असंतृप्त} ठंडी
झिल्ली को उसके नीचे $T_m$ से ऊपर अकड़ाकर, और गरम
झिल्ली को किसी ठंडे दिन जेल बनने से बचाकर।
\textbf{8.} $30\,000\times 38 = \qty{1140}{MJ}$।
\textbf{9.} $1140/60 = 19$ दिन।
\textbf{10.} शुष्क: $\num{1.14e6}/17 = \qty{67}{kg}$; जलयोजित:
\qty{268}{kg}।
\textbf{11.} $268 - 30 = \qty{238}{kg}$, यानी ऊँट के द्रव्यमान का लगभग आधा।
\textbf{12.} समूची त्वचा के नीचे \omterm{def:b1:lipids:triglyceride}{वसा} की एक परत शरीर को कुचालित कर देती
और उसे मरुस्थल में ऊष्मा बहाने से रोक देती; \omterm{def:b1:lipids:triglyceride}{वसा} को एक ही कूबड़ में
सिमटाने से बाकी त्वचा ऊष्मा खोने के लिए खुली रह जाती है।
\textbf{13.} $1.07\times 30 = \qty{32}{kg}$ जल।
\textbf{14.} ऑक्सीजन की $2.0\times 30\,000 = \qty{60000}{L} = \qty{60}{m^3}$।
\textbf{15.} हवा की $60/(0.21\times 0.05) = \qty{5700}{m^3}$।
\textbf{16.} जल की $5700\times 30 = \qty{171}{kg}$।
\textbf{17.} जितना मिलता है उससे पाँच गुना खोया: कूबड़ को ऑक्सीकृत करने
में जल पड़ता है। कूबड़ ऊर्जा का ऐसा भंडार है जो ऊँट को बिना भोजन के चलने
देता है; उसकी जल-किफ़ायत उसके वृक्कों से, निर्जलीकरण की उसकी सहनशीलता से
और अपना तापमान चढ़ने देने की उसकी क्षमता से आती है।
\textbf{18.} $500\times 3.5\times 6 = \qty{10.5}{MJ}$ ऊष्मा के रूप में संचित और रात में छोड़ी जाती है; उसे वाष्पित करके हटाने में प्रति दिन $10\,500/2.4 = \qty{4.4}{kg}$ जल पड़ता।
\textbf{19.} कोई झिल्ली कुछ ही प्रतिशत खिंच सकती है, इसलिए किसी
लाल \omterm{def:b1:cell-unit-of-life:cell}{कोशिका} का आयतन दुगना करने के लिए विश्राम की द्विअवतल आकृति में मुड़ी
हुई बहुत अतिरिक्त झिल्ली चाहिए, और ऐसा कोशिकाकंकाल भी जो उसे बिना फटे
खुलने दे: ऊँट की लाल \omterm{def:b1:cell-unit-of-life:cell}{कोशिकाएँ} अंडाकार हैं, और उनमें प्रति आयतन मनुष्य की
\omterm{def:b1:cell-unit-of-life:cell}{कोशिकाओं} से अधिक झिल्ली है।
\textbf{20.} $2\times 80\times 10^{-12}/0.65\times 10^{-18} = \num{2.5e8}$।
\textbf{21.} $\num{2.5e8}\times\num{8e12}\times 40 = \num{7.9e22}$।
\textbf{22.} $\num{7.9e22}\times 750/\num{6e23} = \qty{99}{g}$।
\textbf{23.} पार्श्व विसरण हर चरण पर शीर्ष को जल में और पूँछों को \omterm{def:b1:lipids:triglyceride}{तेल} में
रखता है, और उसमें कुछ नहीं पड़ता; पलटने के लिए ध्रुवी शीर्ष को
जलविरोधी गूदा पार करना पड़ता है, जो एक बड़ी ऊर्जा-रोक है जिसे तापीय
हलचल हफ़्ते में एक बार पार कर पाती है।
\textbf{24.} दोनों पर्णों के बीच कोई अंतर उतनी ही तेज़ी से क्षीण होता है
जितनी तेज़ी से \omterm{def:b1:lipids:fattyacid}{लिपिड} पलटते हैं; प्रति अणु हफ़्ते में एक बार पर, विशिष्ट
\omterm{def:b1:lipids:fattyacid}{लिपिड} आरपार ले जाते एंजाइमों (फ़्लिपेज़) से बनाई कोई असममिति स्वतःस्फूर्त
पलटने के विरुद्ध अनिश्चित काल तक टिकी रहती है।
\textbf{25.} लगभग \qty{240}{kg} की बचत --- यानी अपने शरीर-द्रव्यमान का
लगभग आधा --- \qty{270}{kg} जलयोजित ग्लाइकोजन के बजाय \qty{30}{kg} \omterm{def:b1:lipids:triglyceride}{वसा}
ढोकर; और कूबड़ कोई जल-भंडार नहीं: उसे ऑक्सीकृत करने में साँस में उससे
पाँच गुना अधिक जल खोता है जितना वह बनाता है।
\end{solution}
